$n=1$のときは不適なので, 以降は$n\geq 2$とする. このとき$2^n - 12$は常に4の倍数であるため, これが25の倍数になるときを調べればよい. 以下, 合同式の法は25とする. $2^{10} = 1024 \equiv -1$より, $2^{20}\equiv 1$である. また, $1\leq r\leq 19$を満たす整数$r$に対して$2^r\not\equiv 1$であることもわかる. したがって, $2^n$を25で割った余りは$n$について周期20で変化する. $2^9 \equiv 12$であるので, $2^n - 12\equiv 0$であるのは$n=20k + 9$ ($k$は0以上の整数) と表されるときである. このとき, $\bmod{11}$では
\[
2^n = 2^{20k + 9} = 1024^{2k} \cdot 512 \equiv 1^k \cdot 6 = 6 \pmod{11}
\]
なので, \emph{$2^n$を$11$で割った余りは$6$であることがわかる. }


$2,4,16,32,1024$は条件を満たし, $8,64,128,256,512$は条件を満たさない. そこで以下では$n\geq 11$の状況を考える.

$2^n$の下3桁に注目し, 百の位を$a$, 十の位を$b$, 一の位を$c$とおく. いま$n\geq 11$より$2^n$は4桁以上であり, 最初の桁は必ず1以上である. よって部分的な桁の和$a+b+c+1$は7以下であるから, $a+b+c\leq 6$でなければならない. また, $2^n$は偶数なので$c$は偶数であり, (8の倍数の判定法から) $100a+10b+c$は8で割り切れる. また, $2^n$が5の倍数でないことから, $c\neq 0$である. これらの条件を満たす$a,b,c$は
\[
(a,b,c) = (0,2,4), (1,1,2), (3,1,2) ,(0,3,2)
\]
となる. それぞれの場合について調べる.
\begin{itemize}
\item $(0,2,4)$のとき, 4以上の整数$e$によって$2^n = 10^{e} + 24$と表せる($e=3$のとき$n=10$に注意). このとき右辺は2で4回割れないが, 左辺は割れるので, 解なし.
\item $(1,1,2)$と$(3,1,2)$のとき, $2^n - 12$は100の倍数なので, (1)より$2^n -12 \equiv 5\pmod{11}$. $2^n - 12$は桁の和が4以下の整数なので, $10^k$という形の整数の4つ以下の和で書くことができる. すなわち, $x_1,x_2,x_3,x_4\in \set{0,1}$ (相異なるとは限らない)と2以上の整数$e_1,e_2,e_3,e_4$が存在して
\[
2^n - 12 = x_1\cdot 10^{e_1} + \dots + x_4\cdot 10^{e_4}
\]
と表せる. $\bmod{11}$を取ると
\[
5\equiv x_1(-1)^{e_1} + \dots + x_4(-1)^{e_4}
\]
であるが, 右辺が取りうる値の範囲は$-4$から$4$までの整数であるため, この合同式が成り立つような$x_i, e_i$を取ることはできない.
 \item $(0,3,2)$のとき, $2^n$の桁和は6以上であるが, $2^n$が3の倍数ではないので桁和は6ではない. よって$2^n$の桁和は7であり, とくに$2^n \equiv 7\equiv 1\pmod{3}$なので$n$は偶数である必要がある. 一方で, $2^n - 32$が$5$の倍数なので$2^n\equiv 2\pmod{5}$が必要であり, このようになるのは$n=4k+1$ ($k$は0以上の整数) と表せるときなので, 偶奇で矛盾する. よってこの場合も解は存在しない.
\end{itemize}
以上より, $2,4,16,32,1024$のみである. よって求める$n$は
\begin{center}
\emph{$n=1,2,4,5,10$}.
\end{center}
